\subsection{The limit operator and the two parity sheets}\label{sec:D03}
We use the uniform comparison with the continuous profile (Section~\ref{sec:P06}) and the gap from Section~\ref{sec:P03}.
Let $v_n=(\mathbf1_R,-\mathbf1_C,0_D,0_A)$ and set
\[
 \begin{gathered}
 P_n=\frac{v_nv_n^T}{2n},\qquad J_n=S_n-I+P_n,\\
 S_n=\operatorname{diag}(0_R,0_C,h_D,h_A)+A^T\operatorname{diag}(g)A,
 \qquad g_c=\frac{\pi_c(1-\pi_c)}{1-4\pi_c}.
 \end{gathered}                                         \tag{D3.1}
\]
The matrix $S_n+P_n$ replaces the unique zero eigenvalue of $S_n$ by one.
Thus $\det(I+J_n)=\operatorname{pdet}S_n$ and
$I+J_n\ge\min(1,\gamma)I$ uniformly.

\paragraph{Diagonal and trace.}
In every family, $h_L+\sum_{c\in L}\pi_c=1$, with $h_R=h_C=0$.
Hence $(S_n)_{LL}=1+\delta_L$, where
$\delta_L=\sum_{c\in L}(g_c-\pi_c)=O(n^{-1})$.
The diagonal part satisfies $\sum_L\delta_L^2=O(n^{-1})$, yet
\[
 \operatorname{tr}J_n
 =1+4\sum_c(g_c-\pi_c)
 =1+12\sum_c\frac{\pi_c^2}{1-4\pi_c}
 \longrightarrow1+12I_2.                          \tag{D3.2}
\]
Here we use $\sum\pi_c^3=O(n^{-1})$ and $\sum\pi_c^2\to I_2$.
The trace cannot be obtained by taking the trace of a Hilbert--Schmidt limit.

\paragraph{A fixed space and isometries.}
With the intervals $I_F$ from Section~\ref{sec:preliminares}, let
\[
 \mathcal H=\bigoplus_{F=R,C,D,A}
          L^2(I_F\times\{0,1\},dt/2).
\]
Each family has two sheets. The sheet of a line is determined by the parity of its absolute index,
$i,j,k,l$ in the respective families. At square $(i,j)$, the diagonal indices $i-j$ and $i+j-1$
differ by $2j-1$, so they always have opposite parity. Merging the two parity
classes before constructing the operator would lose this information about which diagonals
can meet at a square.
For each family and parity, with $m_{F,p}$ indices, divide $I_F$
into $m_{F,p}$ equal intervals, called cells, assigned to the indices in increasing order.
The measure of each cell in $\mathcal H$ is
$\mu_{F,p}=|I_F|/(2m_{F,p})=n^{-1}(1+O(n^{-1}))$.
The isometry from the finite-dimensional Euclidean space is
\[
 E_ne_L=\mu_{F,p}^{-1/2}\mathbf1_{\text{cell of }L}.
                                                        \tag{D3.3}
\]
Each cell lies within distance $O(n^{-1})$ of the physical coordinate
of its index: $(i-1/2)/n$, $(j-1/2)/n$, $k/n$, or $l/n$.
This also holds when the two sheets contain different numbers of indices.

\paragraph{Limit kernels between distinct families.}
Define $\mathcal K$ by the following kernels. The blocks in the reverse order are
their adjoints, and the blocks coupling a family to itself are zero:
\[
\begin{array}{ll}
 RC:&\nu(r,c),\\
 RD:&\nu(r,r-d)\mathbf1_{0<r-d<1},\\
 RA:&\nu(r,a-r)\mathbf1_{0<a-r<1},\\
 CD:&\nu(c+d,c)\mathbf1_{0<c+d<1},\\
 CA:&\nu(a-c,c)\mathbf1_{0<a-c<1}.
\end{array}                                             \tag{D3.4}
\]
These five kernels are independent of both sheet indices.
By contrast,
\[
 \begin{gathered}
 \mathcal K_{DA}((d,p),(a,q))
 =\mathbf1_{p\ne q}\,
   \nu\left(\frac{a+d}{2},\frac{a-d}{2}\right)\mathbf1_\Omega(d,a),\\
 \Omega=\left\{(d,a):0<\frac{a+d}{2}<1,\ 0<\frac{a-d}{2}<1\right\}.
 \end{gathered}                                         \tag{D3.5}
\]
Indeed, $k=i-j$, $l=i+j-1$, and $k+l=2i-1$ is odd.
This condition is independent of $n\bmod2$.
The definition of $\Omega$ requires both arguments to lie between zero and one.

\begin{proposicion}[Hilbert--Schmidt convergence]
Let $f_0$ equal $1/\sqrt2$ on both R sheets, $-1/\sqrt2$ on
both C sheets, and zero on the D and A families. Then $\|f_0\|=1$.
With $P_0=|f_0\rangle\langle f_0|$, we have
\[
 \widehat J_n:=E_nJ_nE_n^*\longrightarrow
             J:=\mathcal K+P_0\quad\text{in Hilbert--Schmidt norm}.       \tag{D3.6}
\]
\end{proposicion}
\begin{proof}
The kernel corresponding to an entry for an intersection is
$g_{ij}/\sqrt{\mu_{F,p}\mu_{G,q}}=ng_{ij}(1+O(n^{-1}))
=\nu(x_i,x_j)+O(n^{-1})$.
Outside bands of width $O(n^{-1})$ around the finitely many segments
bounding the intersection regions, the indicators agree
and the kernels converge uniformly because the profile is Lipschitz.
Inside the bands, the kernels remain bounded and the area of the bands tends to zero.
Thus each block converges in $L^2$. The parity condition in (D3.5) holds exactly.
The diagonal $\delta$ vanishes in Hilbert--Schmidt norm.
Finally, $f_n=E_nv_n/\sqrt{2n}\to f_0$ in $\mathcal H$:
its value on an R sheet is $\sqrt{m_{R,p}/n}\to1/\sqrt2$,
and similarly on a C sheet. Hence the corresponding projections converge in Hilbert--Schmidt norm.
\end{proof}

\paragraph{An equivalent description of parity.}
The change of variables $s=(f_{F,0}+f_{F,1})/2$, $o=(f_{F,0}-f_{F,1})/2$
defines a unitary map onto $\bigoplus_F(L^2(I_F,dt)\oplus L^2(I_F,dt))$.
The symmetric sector has four families, and the block between D and A is
$T(d,a)=\tfrac12\nu((a+d)/2,(a-d)/2)\mathbf1_\Omega$.
In the antisymmetric sector, only the block
$\left(\begin{smallmatrix}0&-T\\-T^*&0\end{smallmatrix}\right)$ on the D and A families remains.
Averaging over parity before constructing the limit would remove this block.
